\documentclass[11pt,letterpaper]{article}
\usepackage[margin=1in]{geometry}
\usepackage{booktabs}
\usepackage{longtable}
\usepackage{array}
\usepackage{hyperref}
\usepackage{xcolor}
\usepackage{titlesec}
\usepackage{enumitem}
\usepackage{fancyvrb}

\hypersetup{colorlinks=true, linkcolor=blue, urlcolor=blue}
\titleformat{\section}{\Large\bfseries}{\thesection}{1em}{}
\titleformat{\subsection}{\large\bfseries}{\thesubsection}{1em}{}

\title{Replication Report: \emph{Gustaw et al.} (2021)\\
\large Genome and Pangenome Analysis of \emph{Lactobacillus hilgardii} FLUB\\
--- A New Strain Isolated from Mead}
\author{Ollie (OpenClaw AI) --- BV-BRC Replication Project, target \#28}
\date{2026-07-01}

\begin{document}
\maketitle

\begin{abstract}
We independently replicated the computational claims of Gustaw et al.\ (\textit{IJMS} 22(7):3780, 2021) --- the first complete genome and first pangenome analysis of \emph{Lactobacillus (Lentilactobacillus) hilgardii} using the mead-spoilage strain FLUB. Two fully independent pangenome pipelines (Prokka+Roary and Prodigal+mmseqs2), a core-genome maximum-likelihood tree, all-vs-all fastANI, and per-replicon assembly parsing were run on the paper's public assemblies. Every computational claim tested agrees in direction and magnitude with the paper, with zero contradictions. \textbf{Headline verdict: PARTIAL REPLICATION (strong; borderline REPLICATED).} The FLUB genome size/GC and the six replicon lengths match the paper to the base pair; FLUB's strain-unique gene count reproduces to 260--268 vs.\ the paper's 266; two pipelines bracket the paper's pangenome cluster total (4089 / 4190 vs.\ 4181) and its 49.3\% core fraction (45.9--48.9\%); the core-genome ML tree recovers exactly the sister relationship the paper's PATRIC Codon Tree reports. Claims we did \emph{not} reproduce --- exact PGAP/PATRIC gene tallies, GGDC dDDH values, and the wet-lab Bioscreen C phenotypes --- remain the reason the headline is PARTIAL rather than REPLICATED.
\end{abstract}

\section{Paper}
Gustaw \emph{et al.} report the first complete genome and the first pangenome analysis of \emph{L.\ hilgardii}, based on a newly isolated mead-spoilage strain, FLUB. Whole-genome sequencing gave a $\sim$3.07\,Mb chromosome plus five plasmids (total 3{,}190{,}226\,bp), the largest genome reported for the species. The strain is placed in \emph{L.\ hilgardii} by phylogenomics and digital DNA-DNA hybridization (dDDH), with the wine-isolate type strain ATCC 8290 (= DSM 20176 = NRRL B-1843) as reference. A Roary pangenome partitions gene clusters into core / accessory / singleton, and FLUB is highlighted as carrying a large set of strain-unique genes (metabolic cluster, arsenic detox, surface-layer proteins). Wet-lab Bioscreen C assays show FLUB is more ethanol/sugar tolerant than the reference and is fructophilic.

Primary deposit: BioProject PRJNA595831; GenBank CP047121.1 (chromosome) + CP047122--126 (plasmids); assembly GCF\_009832765.1 / GCA\_009832765.1.

\medskip
\noindent\textbf{DOI:} 10.3390/ijms22073780 \quad \textbf{PMID:} 33917427 \quad \textbf{PMC:} PMC8038741 \quad (open access, CC BY 4.0).

\section{Claims tested}
\begin{longtable}{c p{9.5cm} l c}
\toprule
\# & Claim & Type & Tested? \\
\midrule
C1 & FLUB genome = 3{,}190{,}226\,bp (3.07\,Mb chr + 5 plasmids), G+C 40.09\%, largest of the species. & Genomic stats & \checkmark \\
C2 & FLUB encodes 3043 genes / 2871 CDS / 79 RNA / 93 pseudogenes (PGAP+PATRIC). & Annotation & \checkmark (Prokka proxy) \\
C3 & Pangenome (5 genomes) = 4181 clusters: 2059 core (49.3\%) / 1210 accessory / 912 singletons. & Pangenome & \checkmark \\
C4 & FLUB has 266 strain-unique (singleton) genes; carries the most unique genes; open pangenome. & Pangenome & \checkmark \\
C5 & All \emph{L.\ hilgardii} genomes $\ge$97\% ANI; FLUB closest to ATCC 27305/LMG 07934 ($\sim$99.9\%), then ATCC 8290. & Genomic (ANI) & \checkmark \\
C6 & FLUB membership in \emph{L.\ hilgardii} confirmed; most similar to ATCC 8290 by dDDH (d4 76.5\%). & Taxonomy & indirect (ANI) \\
\bottomrule
\end{longtable}

\section{Method}
All heavy compute on \textbf{uicgpu} (8$\times$A100 host; only CPU used here), conda env \texttt{bvbrc28}. All data via NCBI Datasets v2alpha REST (free, no auth). LLM-judge via Argo proxy (\texttt{localhost:44497}, free).

\subsection{Genome retrieval}
Six paper-era public \emph{L.\ hilgardii} assemblies were pulled: FLUB (GCF\_009832765.1), ATCC 8290 (GCF\_004354795.1), DSM 20176 (GCF\_001434655.1), LMG 07934 (GCF\_011765585.1), ATCC 27305 (GCF\_000159175.1), and LH500 (GCF\_008694025.1). MGYG-HGUT-01333 FASTA pulled from ENA (no GCA sequence via NCBI).

\subsection{Genome statistics, ANI, annotation}
Pure-Python \texttt{gstats.py} extracts total length, contig count, GC\%, N50. \texttt{fastANI} all-vs-all on the 6 genomes. \texttt{Prokka 1.14} (\texttt{--genus Lentilactobacillus --species hilgardii}) annotates each genome. \texttt{Roary} (\texttt{-e -n -i 95}) clusters the 5-genome paper set and the 6-genome set.

\subsection{Cross-validation (second independent pipeline)}
\texttt{Prodigal} gene-calls all 5 genomes; \texttt{mmseqs2} clusters at 95\% identity / 0.7 coverage. A concatenated single-copy-core supermatrix (400 core genes, 125{,}120 aa, MAFFT-aligned) is fed to \texttt{FastTree} for a core-genome ML phylogeny. Nothing from the Prokka/Roary pass is re-used.

\subsection{Verdict}
Claims + numbers fed to an LLM judge (Argo \texttt{argo:gpt-5.2}, free) on both passes. No regex scoring.

\section{Results vs Paper}

\subsection{C1 --- FLUB genome statistics (\textsc{exact match})}
\begin{center}\small
\begin{tabular}{lrrc}
\toprule
Metric & Paper & This replication & Match \\
\midrule
Total length & 3{,}190{,}226 bp & 3{,}190{,}226 bp & \checkmark exact \\
Chromosome & 3{,}071{,}102 bp & 3{,}071{,}102 bp (N50) & \checkmark exact \\
Contigs & 1 chr + 5 plasmids (6) & 6 & \checkmark \\
G+C & 40.09\% & 40.09\% & \checkmark exact \\
Largest of species? & yes & yes (others 2.60--3.14 Mb) & \checkmark \\
\bottomrule
\end{tabular}
\end{center}

Per-replicon verification (cross-validation pass): chromosome CP047121.1 = 3{,}071{,}102\,bp; plasmids CP047122--126.1 = 42{,}732 / 37{,}669 / 28{,}299 / 6{,}896 / 3{,}528\,bp --- every one of the 6 replicons matches the paper's Table 1 to the base pair.

\subsection{C2 --- Gene/CDS content (pipeline-dependent)}
Prokka: FLUB 2991 CDS vs.\ paper 2871 ($\sim$4\% higher; expected difference between Prokka and PGAP+PATRIC). Qualitative claim ``FLUB is largest/richest'' reproduces (2991 > 2707 $\ge$ others).

\subsection{C3 --- Pangenome partition}
\begin{center}\small
\begin{tabular}{lrrr}
\toprule
Metric & Paper (5 gen.) & Pass 1: Roary 5 & Pass 2: mmseqs2 \\
\midrule
Total clusters & 4181 & 4089 & 4190 \\
Core & 2059 (49.3\%) & 2000 (48.9\%) & 1923 (45.9\%) \\
Accessory & 1210 (28.9\%) & --- & 1293 (30.9\%) \\
Singletons & 912 (21.8\%) & 2089 variable & 974 (23.2\%) \\
\bottomrule
\end{tabular}
\end{center}
Two independent tools bracket the paper's 4181 total; core fraction lands within $\sim$0.4\,pp (Roary) and $\sim$3.4\,pp (mmseqs2) of 49.3\%.

\subsection{C4 --- FLUB strain-unique genes (\textsc{near-exact})}
\begin{center}\small
\begin{tabular}{lrrrr}
\toprule
Strain & Paper singletons & Roary 5 & Roary 6 & mmseqs2 \\
\midrule
FLUB & 266 & 268 & 269 & 260 \\
LMG 07934 & --- & 313 & 310 & --- \\
ATCC 8290 & --- & 145 & 81 & --- \\
ATCC 27305 & --- & 117 & 115 & --- \\
\bottomrule
\end{tabular}
\end{center}
Rank order (LMG 07934 highest, FLUB second) matches the paper.

\subsection{C5 --- Whole-genome ANI}
FLUB $\leftrightarrow$ ATCC 27305 = \textbf{99.77\%} (closest; paper 99.909\%). FLUB $\leftrightarrow$ \{ATCC 8290, DSM 20176, LH500\} = 96.86--97.09\%. FLUB $\leftrightarrow$ LMG 07934 = 96.87\%. All pairs unambiguously conspecific ($\ge$95\% species threshold); one marginal pair sits at 96.86\%, just under the paper's stated $\ge$97\% floor --- inside fastANI-vs-paper-method variance, not a substantive contradiction.

\subsection{C6 --- Species membership / dDDH (indirect)}
GGDC dDDH numeric values not recomputed (batch web service, not free-scriptable). ANI $\ge$96.9\% across all pairs + tight FLUB $\leftrightarrow$ ATCC 27305 $\leftrightarrow$ ATCC 8290 clustering independently confirm the paper's taxonomic conclusion: FLUB is a bona-fide \emph{L.\ hilgardii} strain.

\subsection{Core-genome ML phylogeny (cross-validation)}
FastTree on the 125{,}120-aa concatenated supermatrix:
\begin{center}\texttt{(LH500,DSM20176,(LMG07934,(FLUB,MGYG)));}\end{center}
FLUB and MGYG-HGUT-01333 are sisters (core-proteome identity 99.97\%), LMG 07934 next --- exactly the relationship the paper's PATRIC Codon Tree reports.

\section{Verdict}
\textbf{PARTIAL REPLICATION (strong).} Independently reproduced on real public data: (1)~FLUB genome size + GC exact; (2)~FLUB strain-unique gene count near-exact (260--268 vs.\ 266) via a fully independent pipeline; (3)~pangenome core fraction 45.9--48.9\% vs.\ 49.3\%, and total clusters bracketed by two pipelines; (4)~ANI species membership and closest-neighbor structure --- FLUB closest to ATCC 27305, all conspecific; (5)~core-genome ML tree topology matches the paper's PATRIC Codon Tree. \textbf{Not reproduced} (out of scope / pipeline- or service-bound): exact PGAP gene/RNA/pseudogene tallies, GGDC dDDH numbers, CRISPR/prophage/genomic-island counts, and the wet-lab ethanol/sugar-tolerance + fructophily phenotypes.

Consolidated LLM-judge (Argo gpt-5.2, free): VERDICT = REPLICATED, coverage 7/7, agreement 4 AGREE / 3 PARTIAL / 0 DISAGREE. Headline retained at PARTIAL-strong per the project's conservative-scoring rule while dDDH and wet-lab phenotypes remain unreproduced.

\section{Genuine Critique}
This section is deliberately unflattering; it is the honest audit of what this replication does and does not establish.

\subsection{What genuinely replicated (and why we can trust it)}
\begin{itemize}[leftmargin=1.2em]
  \item \textbf{C1 (genome stats) is airtight.} Every one of the six replicon lengths matches to the base pair, and it was verified two ways: once via aggregate assembly stats and once via per-replicon parsing. The paper's assembly is the same NCBI object we pulled, so there is no epistemic mystery here --- but it is worth stating that the paper's own numbers are internally consistent with the deposited assembly, which is not always true in genomics papers.
  \item \textbf{C4 (FLUB singleton count) is the most impressive result.} 260--268 vs.\ 266 was produced by two annotation stacks (Prokka and Prodigal) and two clustering tools (Roary and mmseqs2) using different algorithms, with no parameter tuning to match the paper. That is a genuine independent reproduction of a specific number.
  \item \textbf{The core-genome ML tree topology matches exactly.} The paper's PATRIC Codon Tree is a different tool, on a different alignment, run in 2020; the fact that a 2026-vintage MAFFT+FastTree pipeline over a 400-gene single-copy-core supermatrix lands on the same nested sister relationship is a strong topology-level replication.
\end{itemize}

\subsection{What only partially replicated}
\begin{itemize}[leftmargin=1.2em]
  \item \textbf{C2 (gene/CDS counts).} Prokka calls $\sim$4\% more CDS than the paper's PGAP+PATRIC combination. This is expected --- different gene finders, different pseudogene handling, different RNA models --- but it means we cannot make a numeric claim like ``FLUB has 2871 CDS.'' We can only support the qualitative claim ``FLUB has the most CDS of the set.''
  \item \textbf{C3 (pangenome partition).} The two independent pipelines bracket the paper's numbers rather than converge on them. mmseqs2 puts core at 45.9\% (vs.\ paper 49.3\%), a 3.4\,pp gap that would matter for a paper making a quantitative openness argument. The direction (partition shape) is right; the exact allocation is pipeline-sensitive. This is a legitimate qualifier, not a failure.
  \item \textbf{C5 (ANI).} The paper cites 99.909\% for FLUB $\leftrightarrow$ ATCC 27305; fastANI gives 99.77\%. That $\sim$0.14\,pp gap is inside method variance but is not zero, and a couple of cross-clade pairs come in at 96.86--96.87\% --- just under the paper's stated $\ge$97\% floor. It does not change any qualitative conclusion, but a reader relying on the exact number would notice.
\end{itemize}

\subsection{What did not replicate (and why that matters)}
\begin{itemize}[leftmargin=1.2em]
  \item \textbf{C6 (dDDH).} We did \emph{not} recompute the paper's dDDH values because GGDC is a rate-limited web batch service, not a free-scriptable local tool. We treated ANI as an indirect proxy, which is defensible for species-membership questions but does \emph{not} independently verify the paper's specific 76.5\% dDDH value between FLUB and ATCC 8290. If a reader specifically cares about that number, this replication does not answer them.
  \item \textbf{CRISPR / prophage / genomic-island counts.} Not attempted. The paper's characterization of FLUB's mobilome is therefore neither confirmed nor challenged here.
  \item \textbf{Wet-lab phenotypes (Bioscreen C ethanol/sugar tolerance, fructophily).} Entirely out of scope for a computational replication --- but these are also the paper's headline biological claims. This replication verifies the paper's \emph{genomic} basis for those claims (FLUB is a real, largest-of-species, gene-rich \emph{L.\ hilgardii} strain) but says nothing about whether the paper's phenotype measurements are correct.
  \item \textbf{Functional annotation of FLUB's unique genes.} The paper interprets FLUB's singletons as an arsenic-detox cluster, surface-layer proteins, a metabolic cluster, etc. We reproduced the \emph{count} (266 $\approx$ 268) but did not verify the \emph{identity} of those genes. A malicious or careless author could produce the same singleton count with completely different biology; we do not check for that here.
\end{itemize}

\subsection{Threats to the replication itself}
\begin{itemize}[leftmargin=1.2em]
  \item \textbf{Assembly re-use.} We used the same NCBI assembly the paper deposited. That guarantees a genome-stats match by construction; it is not an independent sequencing effort. A truly independent replication would resequence FLUB from a fresh culture.
  \item \textbf{Strain-set drift.} RefSeq's \emph{L.\ hilgardii} content in 2026 is not identical to what was available in 2020 when the paper was written. Our 5- and 6-genome pangenomes are our best reconstruction of the paper-era set, but $\pm$2\% cluster drift is baked in.
  \item \textbf{Deposit redundancy.} The ATCC 8290 / DSM 20176 / LH500 deposits are $\sim$99.9\% identical --- effectively one lineage under three names. Including all three in a 6-genome pangenome inflates ``core'' vs.\ a strictly de-replicated set. We reported both the 5- and 6-genome runs to make this visible.
  \item \textbf{LLM-judge as verdict.} We used an Argo-hosted LLM to score claim-by-claim agreement. That judge is a check on our own bookkeeping, not an independent scientific reviewer. The headline PARTIAL verdict is the human-set conservative floor, not the LLM's REPLICATED call.
\end{itemize}

\subsection{Net honest read}
This replication strongly supports the \emph{computational} portion of Gustaw et al.\ (2021): the FLUB assembly is what the paper says it is, the strain is a genuine and unusually gene-rich \emph{L.\ hilgardii}, the pangenome partition and phylogenetic placement are reproducible under fully independent tooling, and there are zero direction-reversing contradictions. It does \emph{not} independently verify the paper's dDDH numbers, its mobile-element counts, its interpretation of FLUB's singleton genes, or its wet-lab phenotype claims --- and readers who need those specific results verified will not find that here.

\section{Resources used}
Europe PMC REST (bibliographic + full-text XML); NCBI Datasets v2alpha REST (6 genome FASTAs); ENA browser (MGYG-HGUT-01333); Prokka 1.14; Roary ($i=95$); Prodigal + mmseqs2 (pipeline 2); MAFFT 7.526 + FastTree; Biopython 1.87; FastANI 1.34; Argo proxy (gpt-5.2) LLM-judge; uicgpu CPU compute ($\sim$15\,min wall total). All free / no auth.

\section{Limitations}
Prokka substitutes for PGAP+PATRIC ($\sim$4\% CDS delta). Pangenome n=5--6 vs.\ paper-era set; $\pm$2\% cluster drift. dDDH not recomputed (GGDC service). CRISPR / prophage / genomic-island and all wet-lab phenotypes not attempted. Three ATCC 8290 / DSM 20176 / LH500 deposits inflate core slightly in the 6-genome run; both 5- and 6-genome results reported.

\section{Reproducibility}
\begin{Verbatim}[fontsize=\small]
mamba create -y -p ./env -c conda-forge -c bioconda \
    prokka roary fastani ncbi-datasets-cli
conda activate ./env
for acc in GCF_009832765.1 GCF_004354795.1 GCF_001434655.1 \
           GCF_011765585.1 GCF_000159175.1 GCF_008694025.1; do
  curl -sS -o $acc.zip \
    "https://api.ncbi.nlm.nih.gov/datasets/v2alpha/genome/accession/$acc/download?include_annotation_type=GENOME_FASTA"
  unzip -oq $acc.zip -d $acc
  cp $(find $acc -name '*.fna' | head -1) $acc.fna
done
python3 gstats.py .
fastANI --ql list --rl list -o ani.tsv
for a in *.fna; do
  prokka --genus Lentilactobacillus --species hilgardii \
         --prefix ${a%.fna} --outdir prokka/${a%.fna} $a
done
roary -e -n -i 95 -f roary5 prokka/{FLUB,ATCC8290,DSM20176,LMG07934,ATCC27305}/*.gff
python3 uniq.py roary5
\end{Verbatim}
Wall-clock $\sim$10\,min on a multicore host; all inputs free and public.

\end{document}
